% !TeX root = ../main.tex
\chapter{Przegląd technologii}
\label{chap:technology}

% FIRST DRAFT. Product descriptions are deliberately short and cite the
% vendors' documentation; extend the comparison with books where the
% reviewer expects literature (docs/review-criteria.md, section 2.1).

System opisany w pracy składa się z pięciu warstw: pozyskiwania danych, przetwarzania strumieniowego, uczenia maszynowego, hurtowni z warstwą raportową oraz infrastruktury, na której to wszystko działa.
Dla każdej warstwy istnieje kilka dojrzałych rozwiązań.
W tym rozdziale porównujemy najważniejsze z nich, a na końcu zestawiamy narzędzia, które wybraliśmy, z alternatywami, które odrzuciliśmy.
Ramą dla całego przeglądu jest wymóg tematu pracy: system ma działać na platformie Google Cloud, z wykorzystaniem usług zarządzanych.
Tam, gdzie ten wymóg przesądzał o wyborze, piszemy o tym wprost.

\section{Przetwarzanie strumieniowe}
\label{sec:tech-streaming}

Przetwarzanie wsadowe zakłada, że zbiór danych jest skończony i znany w chwili rozpoczęcia obliczeń.
Przetwarzanie strumieniowe zakłada coś przeciwnego: dane napływają bez końca, a wynik trzeba podawać, zanim strumień się skończy~\cite{kleppmann2017designing}.
Akidau i współautorzy zwrócili uwagę, że przy takich danych trzeba rozróżniać dwa zegary: czas zdarzenia, czyli chwilę, w której coś się wydarzyło, oraz czas przetwarzania, czyli chwilę, w której system się o tym dowiedział~\cite{akidau2015dataflow}.
Opinie napisane w tej samej minucie mogą dotrzeć do systemu w dowolnej kolejności, a niektóre z dużym opóźnieniem.
W tym samym modelu znacznik wodny (ang. \textit{watermark}) szacuje chwilę, do której system zna już wszystkie zdarzenia, i decyduje, kiedy wynik dla okna czasowego można uznać za gotowy.

W naszym projekcie rozróżnienie zegarów ma wymiar bardzo praktyczny.
Odtwarzamy archiwum, w którym półtora roku opinii mieści się w kilku godzinach, więc czas przetwarzania nie mówi nic o tym, kiedy opinia powstała.
Potok musi pobierać czas zdarzenia z samego rekordu.

\subsection{Silniki przetwarzania}
\label{sec:tech-engines}

Model Akidaua i współautorów powstał w Google i stał się podstawą otwartej biblioteki Apache Beam.
Beam oddziela definicję potoku od silnika, który go wykonuje: ten sam program w języku Java, Python lub Go może działać lokalnie, na silniku Apache Flink, Apache Spark albo w usłudze Google Cloud Dataflow~\cite{docbeamguide}.
Dataflow jest usługą zarządzaną, która sama tworzy maszyny robocze, skaluje ich liczbę i przechowuje stan potoku~\cite{docdataflow}.
Włączyliśmy tryb Streaming Engine, w którym stan i przetasowywanie danych trzyma usługa, a nie maszyny robocze.
Dzięki temu każdej maszynie wystarcza dysk 30~GB zamiast domyślnych 400~GB.
Potok pakujemy jako szablon Flex Template, czyli obraz kontenera z kodem i zależnościami, uruchamiany jednym poleceniem z parametrami.

Apache Flink od początku projektowano jako silnik strumieniowy, w którym przetwarzanie wsadowe jest szczególnym przypadkiem przetwarzania strumienia~\cite{carbone2015apache}.
Apache Spark podchodzi do tego od drugiej strony.
Dzieli strumień na małe porcje wsadowe i przetwarza je tym samym silnikiem, którego używa do analizy wsadowej i uczenia maszynowego~\cite{zaharia2016apache}.

Dla tej pracy oba silniki mają tę samą wadę: wymagają utrzymywania klastra albo korzystania z usługi zewnętrznej, a nasz budżet i wymóg pracy na usługach zarządzanych GCP przemawiały za Dataflow.
Wybierając Beam, nie zamykaliśmy sobie przy tym drogi do Flinka czy Sparka, bo potok w Beam można by w razie potrzeby uruchomić na nich bez przepisywania.
Cenę tego wyboru poznaliśmy w trakcie implementacji.
Część transformacji Beam w języku Python, na przykład zapis do BigQuery przez interfejs Storage Write API, jest w rzeczywistości kodem w języku Java uruchamianym obok, co komplikuje obraz kontenera (podrozdział~\ref{sec:impl-pipeline}).

\subsection{Kolejki komunikatów}
\label{sec:tech-messaging}

Między źródłem danych a potokiem stoi kolejka komunikatów, która przyjmuje zdarzenia od wydawców i dostarcza je subskrybentom.
Apache Kafka przechowuje komunikaty w dzienniku podzielonym na partycje i gwarantuje kolejność w obrębie partycji~\cite{wang2015building}.
Google Cloud Pub/Sub jest usługą zarządzaną bez partycji widocznych dla użytkownika; dostarcza komunikaty co najmniej raz, a kolejność gwarantuje wyłącznie dla komunikatów z tym samym kluczem porządkującym, gdy subskrypcja ma włączone porządkowanie~\cite{docpubsub,docpubsubordering}.

Żadna z tych kolejek nie zapewnia globalnej kolejności wszystkich komunikatów bez dodatkowych zabiegów.
W Kafce wymagałoby to jednej partycji, w Pub/Sub jednego klucza porządkującego.
Wybraliśmy drugą możliwość i to rozwiązanie ma konsekwencje dla wydajności, które mierzymy w rozdziale~\ref{chap:evaluation}.
Za Pub/Sub przemawiało to, że nie wymaga utrzymywania brokerów, jest rozliczany za przesłane dane i współpracuje z Dataflow bez dodatkowej konfiguracji.

\section{MLOps i zarządzanie cyklem życia modeli}
\label{sec:tech-mlops}

Sculley i współautorzy z Google opisali system uczący się jako niewielki fragment kodu modelu otoczony dużą ilością kodu do zbierania danych, konfiguracji, monitorowania i serwowania, w którym łatwo zaciągnąć dług techniczny~\cite{sculley2015hidden}.
Z takich obserwacji wyrosło pojęcie \ac{mlops}.
Kreuzberger, Kühl i Hirschl definiują je jako zestaw praktyk, zasad i narzędzi łączących uczenie maszynowe z ciągłą integracją i wdrażaniem, obejmujący automatyzację potoków, wersjonowanie danych, kodu i modeli, śledzenie eksperymentów, rejestr modeli i monitorowanie~\cite{kreuzberger2023machine}.

Narzędziem, które najczęściej pojawia się w tym kontekście, jest MLflow.
Jego twórcy z firmy Databricks opisali je jako otwartą platformę, która rejestruje parametry, metryki i artefakty przebiegów trenowania, pakuje modele w jednolitym formacie i prowadzi ich rejestr~\cite{zaharia2018accelerating}.
Przebiegi można zapisywać w lokalnym katalogu albo na serwerze śledzenia z bazą danych i magazynem artefaktów.
Alternatywą jest Kubeflow, platforma potoków uczenia maszynowego działająca na Kubernetesie, która daje więcej możliwości, ale wymaga klastra.
W ekosystemie Google eksperymenty śledzi też usługa zarządzana Agent Platform Experiments.

Wybraliśmy MLflow, bo jest niezależny od dostawcy chmury, wymaga najmniej infrastruktury i jest wprost wymieniony w temacie pracy jako platforma do zarządzania cyklem życia modelu.
Do trenowania i serwowania użyliśmy natomiast usług Agent Platform: środowiska notatników Workbench, rejestru modeli i punktu końcowego, który udostępnia model przez sieć~\cite{docvertexdeploy}.
Usługi te nosiły wcześniej nazwę Vertex AI i od niej pochodzi identyfikator \code{vertex}, którego używamy w kodzie projektu.

\section{Usługi chmurowe do analityki danych}
\label{sec:tech-cloud}

Chmura obliczeniowa w ujęciu NIST to model, w którym zasoby obliczeniowe udostępnia się na żądanie, przez sieć, z rozliczeniem za faktyczne zużycie i z możliwością szybkiego zwiększania i zmniejszania skali~\cite{mell2011nist}.
Dla nas najważniejsze jest rozliczanie za zużycie, bo od niego zależy, ile kosztuje zasób, który przez większość czasu nic nie robi.
Najdroższe zasoby, czyli zadanie Dataflow i punkt końcowy modelu, istnieją tylko na czas przetwarzania, a reszta usług jest rozliczana za zużycie.

Hurtownią danych w GCP jest BigQuery, kolumnowa baza odpytywana językiem \ac{sql}, bez serwerów do zarządzania, rozliczana za przetworzone dane albo zarezerwowaną moc obliczeniową~\cite{docbigquery}.
Dla nas ważne są trzy jej cechy.
Polecenie \code{ORDER BY} sortuje cały wynik, co wykorzystujemy do ustalenia globalnej kolejności opinii.
Instrukcja \code{MERGE} pozwala scalać wyniki z deduplikacją po identyfikatorze.
Widoki autoryzowane mogą czytać tabelę źródłową, do której odbiorca widoku nie ma dostępu, co pozwala ukryć dane osobowe przed panelem.

Surowe dane leżą w Cloud Storage, magazynie obiektów rozliczanym za zajmowane miejsce.
Kod pobierający dane i odtwarzający strumień działa jako zadania Cloud Run Jobs, czyli kontenery uruchamiane na żądanie, które kończą się po wykonaniu pracy.
Widoki raportowe definiujemy w Dataform, narzędziu wbudowanym w BigQuery, w którym widoki opisuje się w plikach SQL z testami poprawności danych~\cite{docdataform}.
Podobne możliwości daje dbt, ale wymaga osobnego środowiska uruchomieniowego i przechowywania poświadczeń.

Region przetwarzania, \code{europe-central2} w Warszawie, wybraliśmy ze względu na ochronę danych.
Pseudonimy autorów i treści opinii nie opuszczają Unii Europejskiej, co upraszcza zgodność z ogólnym rozporządzeniem o ochronie danych (\ac{rodo})~\cite{rodo2016}.

\section{Infrastruktura jako kod}
\label{sec:tech-iac}

Infrastruktura jako kod (\ac{iac}) to praktyka, w której zasoby, sieci i uprawnienia opisuje się w plikach tekstowych przechowywanych w repozytorium i stosuje automatycznie, zamiast klikać je w konsoli~\cite{morris2021infrastructure}.
Skoro infrastruktura jest kodem, ma też jego wady: błędy, luki w uprawnieniach i potrzebę przeglądu przed każdą zmianą.

W Terraform firmy HashiCorp zasoby opisuje się deklaratywnie, w języku HCL, a narzędzie porównuje opis z zapisanym stanem i wylicza plan zmian, który można przejrzeć przed zastosowaniem~\cite{docterraform}.
Pulumi pozwala opisywać te same zasoby w językach ogólnego przeznaczenia, takich jak Python czy TypeScript.
AWS CloudFormation działa tylko w chmurze Amazon.

Wybraliśmy Terraform z trzech powodów.
Ma dojrzałego dostawcę dla Google Cloud, plan zmian dobrze nadaje się do przeglądu w zgłoszeniu zmian, a deklaratywny opis łatwo sparametryzować przedrostkiem nazw, dzięki któremu kolejne studium przypadku dostaje własny komplet zasobów.
Zmiany infrastruktury przechodzą przez potok ciągłej integracji i wdrażania (\ac{cicd}) w usłudze Cloud Build, który dla zgłoszenia zmian wylicza plan, a po scaleniu czeka na zatwierdzenie przez człowieka.

\section{Wizualizacja danych}
\label{sec:tech-visualisation}

Panel analityczny ma pozwolić odbiorcy w kilka sekund odczytać, jak zmieniały się nastroje, i nie wprowadzać go w błąd.
Few zaleca, aby panel mieścił się na jednym ekranie, pokazywał liczby w kontekście, na przykład względem punktu odniesienia, i unikał ozdobników, które nie niosą informacji~\cite{few2006information}.
Z tych zasad wynikają nasze wymagania wobec panelu: przy każdym wyniku liczność danych, wspólne kolory klas wydźwięku i zaznaczone daty wydarzeń.

Do budowy panelu rozważaliśmy cztery narzędzia.
Data Studio jest bezpłatnym narzędziem Google z wbudowanym połączeniem z BigQuery, które potrafi wykonywać zapytania z uprawnieniami osoby oglądającej raport~\cite{doclookerstudio}.
Grafana jest nastawiona na szeregi czasowe i monitorowanie systemów, a Power BI to rozbudowane narzędzie biznesowe firmy Microsoft.
Streamlit pozwala zbudować panel jako aplikację w języku Python.
Z tej możliwości skorzystaliśmy przejściowo: aplikacja działała w Cloud Run za usługą Identity-Aware Proxy (\ac{iap}), dopóki 2 listopada 2026~r. nie zastąpił jej raport w Data Studio.

\section{Zestawienie wybranych technologii}
\label{sec:tech-summary}

Tabela~\ref{tab:tech-summary} zestawia wybrane technologie z rozważanymi alternatywami i główną przyczyną wyboru.
Wspólnym mianownikiem wyborów jest to, że prawie wszystkie wybrane usługi są zarządzane przez dostawcę i rozliczane za faktyczne zużycie.
Jedynym zasobem, który działa i kosztuje stale, jest baza danych serwera MLflow; ten wyjątek uzasadniamy w podrozdziale~\ref{sec:impl-nlp}.

\begin{table}[hbp]
  \caption{Wybrane technologie i rozważane alternatywy}
  \label{tab:tech-summary}
  \centering
  \small
  \begin{tabular}{|p{0.16\textwidth}|p{0.19\textwidth}|p{0.19\textwidth}|p{0.31\textwidth}|}
    \hline
    Warstwa                  & Wybór                                   & Alternatywy                  & Przyczyna wyboru                                                                     \\
    \hline
    Kolejka komunikatów      & Pub/Sub                                 & Apache Kafka                 & usługa zarządzana, klucze porządkujące, rozliczenie za dane                          \\
    \hline
    Przetwarzanie strumienia & Apache Beam na Dataflow                 & Apache Flink, Apache Spark   & usługa zarządzana, czas zdarzenia, przenośność potoku między silnikami               \\
    \hline
    Hurtownia danych         & BigQuery                                & --                           & globalne sortowanie, MERGE, widoki autoryzowane, brak serwerów                       \\
    \hline
    Archiwum surowych danych & Cloud Storage                           & --                           & niski koszt przechowywania, ładowanie do BigQuery                                    \\
    \hline
    Uruchamianie usług       & Cloud Run Jobs                          & maszyny wirtualne            & kontener działa tylko na czas zadania                                                \\
    \hline
    Model językowy           & XLM-RoBERTa (dostrojony)                & DistilBERT, metody słownikowe, LLM & wielojęzyczność, koszt, powtarzalność                                            \\
    \hline
    Śledzenie eksperymentów  & MLflow                                  & Kubeflow, Agent Platform Experiments & niezależność od dostawcy, mała infrastruktura, wymóg tematu                       \\
    \hline
    Trenowanie i serwowanie  & Agent Platform (Workbench, rejestr, punkt końcowy) & własne maszyny wirtualne & usługa zarządzana, model poza obrazem potoku                                       \\
    \hline
    Warstwa raportowa        & Dataform                                & dbt, SQL w Terraform         & działa w BigQuery, testy danych przy każdym wdrożeniu                                \\
    \hline
    Panel                    & Data Studio                             & Grafana, Power BI, Streamlit & bezpłatny, połączenie z BigQuery, uprawnienia oglądającego                           \\
    \hline
    Infrastruktura           & Terraform                               & Pulumi, CloudFormation       & dojrzały dostawca GCP, przegląd planu zmian, parametryzacja                          \\
    \hline
    CI/CD                    & Cloud Build                             & --                           & integracja z GitHub i GCP, zatwierdzanie zmian                                       \\
    \hline
  \end{tabular}
\end{table}

Kreska w kolumnie alternatyw oznacza, że wybór wynikał bezpośrednio z wymogu pracy na usługach zarządzanych Google Cloud i nie porównywaliśmy szczegółowo rozwiązań innych dostawców.
\todo[inline]{Recenzent oczekuje uzasadnienia wyborów względem alternatyw (docs/review-criteria.md): rozważyć rozbudowę porównania o rozwiązania AWS i Azure, z literaturą.}
W kolejnym rozdziale przechodzimy od narzędzi do wymagań: co dokładnie system ma robić i dla kogo.
